%% ============================================================
%%  CHAPTER 5 — SOFTWARE ARCHITECTURE
%% ============================================================
\chapter{Software Architecture}
\label{ch:software}

\section{Design Philosophy}

The \VIBE{} framework is designed around three principles: \emph{modularity} (each physics component is independently replaceable), \emph{declarativity} (PDEs are expressed symbolically, not procedurally), and \emph{performance} (all numerical operations are GPU-tensor-native).

\section{Overall Architecture}

\begin{figure}[H]
\centering
\begin{tikzpicture}[
  block/.style={rectangle, draw, rounded corners, fill=blue!10,
    text width=3.5cm, align=center, minimum height=1.0cm, font=\small},
  arrow/.style={-Stealth, thick},
  node distance=1.2cm and 0.5cm
]
\node[block] (user) {User Script\\(\texttt{Experiment.py})};
\node[block, below=of user] (solver) {\texttt{ImplicitBatterySolver}\\(pack orchestrator)};
\node[block, below left=of solver] (physics) {\texttt{BatteryPhysics}\\(per-cell PDE eval)};
\node[block, below right=of solver] (solver2) {\texttt{BasicSolver} /\\
\texttt{AdvancedSolver}};
\node[block, below=of physics] (pipeline) {\texttt{OperatorPDEPipeline}\\(equation registry)};
\node[block, below=of pipeline] (models) {PDE Models:\\Spherical, FVM,\\FDM, Chebyshev};
\node[block, below=of solver2] (pack) {Pack modules:\\Balancing,\\Cooling};
\draw[arrow] (user) -- (solver);
\draw[arrow] (solver) -- (physics);
\draw[arrow] (solver) -- (solver2);
\draw[arrow] (physics) -- (pipeline);
\draw[arrow] (pipeline) -- (models);
\draw[arrow] (solver) -- (pack);
\end{tikzpicture}
\caption{High-level software architecture of the \VIBE{} framework.}
\label{fig:arch_overview}
\end{figure}

\section{Symbolic PDE Expression AST}
\label{sec:symbolic_ast}

The framework provides a symbolic abstract syntax tree (AST) for PDE definition, implemented in \texttt{pde\_framework.py} and \texttt{pde\_sim/symbolic/expressions.py}.

\subsection{Expression Nodes}

All AST nodes inherit from the \texttt{Expression} base class, which overloads Python arithmetic operators:

\begin{lstlisting}[language=Python, caption={Expression class with operator overloading}]
class Expression:
    def __add__(self, other): return Add(self, ensure_expression(other))
    def __mul__(self, other): return Multiply(self, ensure_expression(other))
    def __neg__(self):        return Negate(self)
    # ... (sub, div, pow, radd, rmul, etc.)
\end{lstlisting}

Leaf nodes include:
\begin{itemize}
  \item \texttt{Constant(value)} — a literal scalar
  \item \texttt{Variable(name)} / \texttt{Field(name)} — a solved state (unknown)
  \item \texttt{Parameter(name)} / \texttt{Param(name)} — a runtime coefficient
\end{itemize}

Differential operator nodes include:
\begin{itemize}
  \item \texttt{Grad(expr)} — spatial gradient $\nabla(\cdot)$
  \item \texttt{Div(expr)} — divergence $\nabla \cdot (\cdot)$
  \item \texttt{Laplacian(expr)} — $\nabla^2(\cdot) = \nabla \cdot \nabla(\cdot)$
  \item \texttt{SphericalDiv(expr, radius)} — $(1/r^2)\partial_r(r^2 \cdot)$
  \item \texttt{Dt(expr)} — time derivative $\partial_t(\cdot)$ (for LHS notation)
\end{itemize}

With these primitives, the solid diffusion PDE can be declared as:

\begin{lstlisting}[language=Python, caption={Declarative PDE for solid diffusion}]
c = Field("cs_n")
D = Param("diffusivity")
r = Param("particle_radius")
# Declares: dc/dt = (1/r^2) * d/dr(r^2 * D * dc/dr)
rhs = SphericalDiv(D * Grad(c), r)
\end{lstlisting}

and the electrolyte transport PDE as:

\begin{lstlisting}[language=Python, caption={Declarative PDE for electrolyte transport}]
ce = Field("ce")
D_eff = Param("flux_coefficient")   # negative of diffusivity
source = Param("source")
rhs = -Div(D_eff * Grad(ce)) + source
\end{lstlisting}

\subsection{Expression Evaluator (\texttt{discretize})}

The \texttt{discretize(expr, context)} function recursively walks the AST and replaces each node with its discretised tensor equivalent:

\begin{align}
  \texttt{Grad}(u) &\rightarrow \mathbf{G}_h \cdot \mathbf{u} \label{eq:discretize_grad}\\
  \texttt{Div}(F) &\rightarrow \mathbf{D}_h \cdot \mathbf{F} \label{eq:discretize_div}\\
  \texttt{Multiply}(a, b) &\rightarrow a_h \odot b_h \label{eq:discretize_mul}
\end{align}
where $\mathbf{G}_h$ and $\mathbf{D}_h$ are the discrete gradient and divergence matrices of the active operator backend.

\section{State Layout Registry}
\label{sec:state_layout}

The \texttt{StateLayout} class manages the flat state vector $\mathbf{y} \in \mathbb{R}^{N_\text{state}}$:

\begin{lstlisting}[language=Python, caption={State registration example}]
layout = StateLayout()
layout.register("cs_n",        size=Nr_n,  initial=29866.0, scale=30000.0)
layout.register("cs_p",        size=Nr_p,  initial=17038.0, scale=30000.0)
layout.register("electrolyte", size=Nel,   initial=eps*ce0,  scale=1000.0)
layout.register("Lsei",        size=1,     initial=5e-9,    scale=1e-8)
layout.register("temperature", size=1,     initial=298.15,  scale=300.0)
\end{lstlisting}

Each field occupies a contiguous slice of $\mathbf{y}$. The slices are computed automatically as fields are registered. Key methods:
\begin{itemize}
  \item \texttt{get(y, name)} — returns the tensor slice for field \texttt{name}
  \item \texttt{pack(derivatives)} — assembles the full $\dot{\mathbf{y}}$ vector
  \item \texttt{initial\_state(params\_list)} — builds the initial condition batch
  \item \texttt{scale\_vector()} — returns $\mathbf{S}$ for Newton convergence scaling
  \item \texttt{nonnegative\_mask()} — boolean mask for line-search enforcement
\end{itemize}

\section{Operator PDE Pipeline}
\label{sec:operator_pipeline}

The \texttt{OperatorPDEPipeline} is a registry that maps state names to discretised PDE models. Three model types are supported:

\begin{description}
  \item[\texttt{SphericalParticlePDEModel}] for solid diffusion in spherical coordinates (Chebyshev backend). Manages flux BC at the surface and symmetry BC at the centre.
  \item[\texttt{ConservativeFluxPDEModel}] (alias \texttt{ElectrolytePDEModel}) for conservative $-\nabla \cdot (\mathbf{D}\nabla c) + S$ form with FVM or Chebyshev backends.
  \item[\texttt{GeneralOperatorPDEModel}] for arbitrary expression-based PDEs without special boundary flux handling.
\end{description}

Adding a new PDE requires only:
\begin{enumerate}
  \item Declaring an \texttt{OperatorEquationSpec} with the symbolic RHS expression.
  \item Adding a \texttt{StateEquationSpec} entry to the \texttt{STATE\_EQUATIONS} registry in \texttt{model\_registry.py}.
  \item Providing runtime parameter resolvers (can be context values, parameter lookups, or callables).
\end{enumerate}

\section{Model Registry and State Equations}
\label{sec:model_registry}

The \texttt{STATE\_EQUATIONS} dictionary in \texttt{model\_registry.py} is the single source of truth for all solved states:

\begin{table}[H]
\centering
\caption{Registered state equations in the \VIBE{} framework.}
\label{tab:state_equations}
\small
\begin{tabular}{llcp{5.5cm}}
\toprule
\textbf{Key} & \textbf{Order} & \textbf{Size} & \textbf{Governing Equation} \\
\midrule
\texttt{cs\_n} & 10 & $N_{r,n}$ & $\partial_t c_{s,n} = D_{s,n} r^{-2}\partial_r(r^2 \partial_r c_{s,n})$ \\
\texttt{cs\_p} & 20 & $N_{r,p}$ & $\partial_t c_{s,p} = D_{s,p} r^{-2}\partial_r(r^2 \partial_r c_{s,p})$ \\
\texttt{electrolyte} & 30 & $N_{x,n}+N_{x,s}+N_{x,p}$ & $\partial_t \hat{c}_e = -\nabla\cdot(-D_\text{eff}\nabla c_e) + S_e$ \\
\texttt{stress} & 40 & $N_\text{el}$ & $\partial_t \sigma = -\nabla\cdot(-D_\sigma\nabla\sigma) + \lambda(c_e - c_{e,0}) - \mu\sigma$ (optional) \\
\texttt{sei} & 80 & $N_\text{sei}$ & Reserved (algebraic) \\
\texttt{Lsei} & 90 & 1 & $\dot{L}_\text{SEI} = -(V_\text{bar}/(\mathcal{F} z)) j_\text{SEI}$ \\
\texttt{temperature} & 1000 & 1 & $\dot{T} = (Q_\text{gen} - Q_\text{cool}) / (\rho C_p V)$ \\
\bottomrule
\end{tabular}
\end{table}

The \texttt{order} field controls the position of each state in the flat vector. Temperature is always placed last (order 1000) because the solver appends pack-level conduction to the final state entry.

\section{Discretisation Backends}
\label{sec:disc_backends}

Three operator backends are implemented, all sharing the same interface (\texttt{gradient()}, \texttt{divergence()}, \texttt{face\_coefficients()}):

\begin{description}
  \item[\texttt{ChebyshevRegionOperators}] Block-diagonal Chebyshev differentiation matrices, one per mesh region. Spectral accuracy for smooth profiles.
  \item[\texttt{FiniteVolumeOperators}] Cell-centred FVM with harmonic-mean interface diffusivities. Exact flux conservation. Preferred for through-cell transport.
  \item[\texttt{FiniteDifferenceOperators}] Node-centred central differences. Second-order convergence. Included for comparison studies.
\end{description}

Selection is made at runtime via the \texttt{config} dictionary:
\begin{lstlisting}[language=Python, caption={Runtime discretisation selection}]
config = {
    "electrolyte_spatial_method": "finite_volume",  # or "chebyshev", "finite_difference"
    "solid_spatial_method":        "chebyshev",
}
\end{lstlisting}

\section{Composite Mesh}
\label{sec:composite_mesh}

The \texttt{CompositeMesh} class stitches together region-wise meshes into a single global mesh object. Each region is defined by a length and a node count:

\begin{lstlisting}[language=Python, caption={Composite electrolyte mesh construction}]
region_spec = {
    'negative':  {'length': Ln, 'N': Nx_n},
    'separator': {'length': Ls, 'N': Nx_s},
    'positive':  {'length': Lp, 'N': Nx_p},
}
mesh = CompositeMesh(region_spec, method='finite_volume')
\end{lstlisting}

The mesh stores: node coordinates, cell widths, face positions, and region-to-slice mappings, enabling each operator to act on the correct sub-domain.
